\documentclass[11pt]{article}

\usepackage[margin=1in]{geometry}
\usepackage{enumitem}
\usepackage{hyperref}
\usepackage{booktabs}
\usepackage{array}
\usepackage{tabularx}
\usepackage{cite}

\setlength{\parindent}{0pt}
\setlength{\parskip}{5pt}

\begin{document}

% =========================================================
% Project Information
% =========================================================

\begin{center}
    {\Large \textbf{CS/CMDA/STAT 3654: Final Project Proposal}}\\[8pt]

    {\large \textbf{Project Title: [Your Project Title]}}\\[6pt]

    \textbf{Team Members:} 
    [Member 1], [Member 2], [Member 3]
\end{center}

\vspace{0.5em}

% =========================================================
% PART 1: PROJECT DESCRIPTION
% =========================================================

\section*{1. Project Description}

\subsection*{1.1 Project Topic and Problem Definition}

Introduce your project topic and application domain. Clearly define the problem your project aims to address using data analytics. Briefly describe the motivation for studying this problem in the project.

\textit{[Replace this text with your content.]}

\subsection*{1.2 Dataset(s)}

Identify the dataset(s) you currently plan to use. Briefly describe the data source, what the data contains, and why it is relevant to your project. Include links or references to the data source when available.

\textit{[Replace this text with your content.]}

\subsection*{1.3 Initial Data Analysis}

Provide a brief description and basic statistical analysis of the dataset(s). Depending on your data, this may include:

\begin{itemize}[leftmargin=1.5em, itemsep=0pt, topsep=2pt]
    \item Dataset size and key variables
    \item Summary statistics (e.g., mean, median, minimum, maximum, standard deviation)
    \item Missing values (if any)
    \item Other relevant dataset characteristics
\end{itemize}

\textit{[Replace this text with your content. You may include small tables or figures if appropriate.]}

\subsection*{1.4 Planned Data Analytics Technique(s)}

Describe the data analytics technique(s) you currently plan to use. Briefly explain how the proposed technique(s) will help address your project problem.

\textit{[Replace this text with your content.]}

% Please add your cited papers to reference.bib.
Example citation: \cite{Li_2020_CVPR}. The cited paper will appear in the References section.


% =========================================================
% PART 2: MANAGEMENT PLAN
% =========================================================

% \newpage

\section*{2. Management Plan}

\subsection*{2.1 Team Responsibilities}

Describe the planned responsibilities of each team member.

For this proposal, please also indicate the percentage contribution of each team member and clearly describe how the work is split. Team members are expected to contribute equally to each submission.

\begin{center}
\begin{tabularx}{\textwidth}{
    >{\raggedright\arraybackslash}p{0.20\textwidth}
    >{\centering\arraybackslash}p{0.18\textwidth}
    >{\raggedright\arraybackslash}X}
\toprule
\textbf{Team Member} & \textbf{Contribution} & \textbf{Work Completed for This Proposal} \\
\midrule
\textit{Member 1} & \textit{XX\%} & \textit{Description} \\
\textit{Member 2} & \textit{XX\%} & \textit{Description} \\
\textit{Member 3} & \textit{XX\%} & \textit{Description} \\
\bottomrule
\end{tabularx}
\end{center}

\subsection*{2.2 Project Timeline and Milestones}

Provide a tentative timeline for the major stages of your project, such as data preparation, methodology and code development, experiments and analysis, presentation preparation, and final deliverables.

\begin{center}
\begin{tabularx}{\textwidth}{
    >{\raggedright\arraybackslash}p{0.25\textwidth}
    >{\raggedright\arraybackslash}X}
\toprule
\textbf{Time / Milestone} & \textbf{Planned Work} \\
\midrule
\textit{Date or Week} & \textit{Planned work} \\
\textit{Date or Week} & \textit{Planned work} \\
\textit{Date or Week} & \textit{Planned work} \\
\textit{Date or Week} & \textit{Planned work} \\
\bottomrule
\end{tabularx}
\end{center}


% =========================================================
% REFERENCES
% References do not count toward the two-page limit.
% =========================================================

% \newpage

\bibliographystyle{plain}
\bibliography{reference}

\end{document}
